\section{Дерево палиндромов и суффиксный автомат}

В первой половине сегодняшней лекции мы разберём структуру данных, придуманную в 2014 году уральским школьником Михаилом Рубинштейном, "--- \textbf{дерево палиндромов} (англ. \textit{Eertree} или \textit{palindromic tree}). Во второй половине мы познакомимся с понятием правого контекста, теоремой Майхилла--Нероуда и приступим к изучению \textbf{суффиксного автомата}.

\subsection{Дерево палиндромов (Eertree)}

Напомним, что строка $\alpha$ называется \textbf{палиндромом}, если она читается одинаково слева направо и справа налево ($\alpha = \alpha^R$). Сколько вообще различных подстрок-палиндромов может быть у строки длины $n$? У произвольной строки количество различных подстрок может достигать $\frac{n(n+1)}{2} = \Theta(n^2)$. Однако для палиндромов верна линейная оценка.

\begin{lemma}[О числе различных палиндромов]\label{lem:pal_count}
    В любой строке $s$ длины $n$ содержится не более $n$ попарно различных непустых подстрок-палиндромов.
\end{lemma}

\begin{proof}
    Будем доказывать утверждение индукцией по длине строки $n$. База при $n = 0$ очевидна (непустых подстрок нет). Покажем, что при приписывании одного нового символа $c = s[i]$ в конец строки $s[0 \dots i-1]$ количество различных подстрок-палиндромов увеличивается \textbf{не более чем на единицу}.

    Все новые палиндромы, которые могли появиться после дописывания $s[i]$, обязаны заканчиваться в позиции $i$. Если в позиции $i$ заканчивается не более одного палиндрома, утверждение сразу выполнено. Пусть в позиции $i$ заканчивается несколько палиндромов; выберем среди них самый длинный палиндром $P$, и пусть $P'$ "--- любой другой (более короткий) палиндром, также заканчивающийся в позиции $i$.

    \begin{center}
    \begin{tikzpicture}[scale=0.95]
        \draw[thick] (0, 0) rectangle (9, 0.7);
        \draw[thick, fill=NavyBlue!12] (2, 0) rectangle (9, 0.7);
        \draw[thick, fill=BrickRed!20] (2, 0) rectangle (4.2, 0.7) node[pos=0.5, font=\small] {$(P')^R = P'$};
        \draw[thick, fill=BrickRed!20] (6.8, 0) rectangle (9, 0.7) node[pos=0.5, font=\small] {$P'$};
        \draw[dashed, thick, NavyBlue] (5.5, -0.3) -- (5.5, 1.1);
        \node[above, font=\footnotesize, text=NavyBlue] at (5.5, 1.1) {ось симметрии $P$};
        \draw[decorate, decoration={brace, mirror, amplitude=6pt}, thick] (2, -0.15) -- (9, -0.15) node[midway, below=6pt, font=\small] {наибольший суффикс-палиндром $P$};
        \node[above, font=\footnotesize] at (9, 0.7) {$i$};
    \end{tikzpicture}
    \end{center}

    Так как $P'$ короче $P$ и оба они заканчиваются в позиции $i$, строка $P'$ является собственным суффиксом палиндрома $P$. Отразим суффикс $P'$ относительно центра симметрии палиндрома $P$: поскольку $P$ симметричен, в начале палиндрома $P$ стоит префикс, равный развёрнутой строке $(P')^R$. Но $P'$ сам является палиндромом, поэтому $(P')^R = P'$.

    Значит, строка $P'$ является собственным префиксом палиндрома $P$, а потому её правый конец в этом вхождении лежит строго левее позиции $i$. Следовательно, палиндром $P'$ уже встречался в префиксе $s[0 \dots i-1]$, и единственным кандидатом на роль нового палиндрома является самый длинный суффикс-палиндром $P$.
\end{proof}

\subsubsection{Структура двух корней и суффиксных ссылок}

Раз различных палиндромов не более $n$, хочется хранить каждый из них как вершину некоторого бора. Заметим, что если к палиндрому $\alpha$ приписать одну и ту же букву $c$ одновременно слева и справа, получится новый палиндром $c\alpha c$ длины $|\alpha| + 2$. Обратно, удаляя у палиндрома длины $\ge 2$ крайние буквы слева и справа, мы снова получаем палиндром.

Поэтому в дереве палиндромов ребро $v \xrightarrow{c} u$ с меткой $c \in \Sigma$ означает переход от палиндрома $v$ к палиндрому $u = c + v + c$, то есть фактически в боре хранятся только правые половинки палиндромов.

При таком определении переходов длина строки увеличивается сразу на $2$: $\len(u) = \len(v) + 2$. Из-за этого палиндромы чётной и нечётной длины никогда не смешиваются по рёбрам дерева:
\begin{itemize}
    \item Все палиндромы \textbf{чётной длины} растут из пустой строки $\epsilon$. Создадим для неё вершину $0$ "--- \textbf{чётный корень} с длиной $\len(0) = 0$.
    \item Все палиндромы \textbf{нечётной длины} начинаются с однобуквенных палиндромов $c$ длины $1$. Чтобы однобуквенный палиндром $c$ получался по тому же правилу $\len(u) = \len(v) + 2$, создадим вершину $1$ "--- \textbf{нечётный корень} (<<чёрную дыру>>) с формальной длиной $\len(1) = -1$! Когда мы переходим из нечётного корня по ребру с буквой $c$, одна из двух букв $c$ <<поглощается>> чёрной дырой, и на выходе получается палиндром $c$ длины $-1 + 2 = 1$.
\end{itemize}

Как и в автомате Ахо--Корасик, определим для каждой вершины $v$ \textbf{суффиксную ссылку} $\link(v)$ как вершину, отвечающую самому длинному собственному суффиксу палиндрома $v$, который сам является палиндромом. Для любого непустого палиндрома (при $\len(v) \ge 1$) такой собственный суффикс-палиндром существует (хотя бы пустая строка $0$). Для самого чётного корня $0$ положим по определению $\link(0) = 1$ (ссылка ведёт в нечётный корень).

\begin{center}
\begin{tikzpicture}[
    scale=1.0,
    vertex/.style={circle, draw=black!80, thick, minimum size=9.5mm, inner sep=2.5pt, font=\footnotesize},
    root/.style={vertex, fill=NavyBlue!15},
    treeedge/.style={draw=black!80, -Stealth, thick},
    linkedge/.style={draw=BrickRed, -Stealth, dashed, thick}
]
    \node[root] (1) at (-2.0, 2.6) {$-1$};
    \node[root] (0) at (2.0, 2.6) {$0$};

    \node[vertex] (a) at (-3.4, 0.9) {$\mathtt{a}$};
    \node[vertex] (b) at (-0.7, 0.9) {$\mathtt{b}$};
    \node[vertex] (bb) at (2.0, 0.9) {$\mathtt{bb}$};

    \node[vertex] (aba) at (-0.7, -1.3) {$\mathtt{aba}$};
    \node[vertex] (abba) at (2.0, -1.3) {$\mathtt{abba}$};

    \draw[treeedge] (1) -- node[above left, font=\small\bfseries, text=NavyBlue] {$\mathtt{a}$} (a);
    \draw[treeedge] (1) -- node[left=3pt, pos=0.75, font=\small\bfseries, text=NavyBlue] {$\mathtt{b}$} (b);
    \draw[treeedge] (0) -- node[right, font=\small\bfseries, text=NavyBlue] {$\mathtt{b}$} (bb);
    \draw[treeedge] (b) -- node[right=3pt, pos=0.25, font=\small\bfseries, text=NavyBlue] {$\mathtt{a}$} (aba);
    \draw[treeedge] (bb) -- node[right, font=\small\bfseries, text=NavyBlue] {$\mathtt{a}$} (abba);

    % Suffix links
    \draw[linkedge] (0) -- (1);
    \draw[linkedge, bend left=8] (a) to (0);
    \draw[linkedge, bend right=12] (b) to (0);
    \draw[linkedge] (bb) to (b);
    \draw[linkedge, bend left=14] (aba) to (a);
    \draw[linkedge, bend left=5] (abba) to (a);
\end{tikzpicture}
\par\nopagebreak\medskip
\small\textit{Рис. 4.1. Дерево палиндромов для строки $s = \mathtt{abbaaba}$. Сплошными стрелками показаны переходы $\go(v, c)$, переводящие $v \to c + v + c$, а красным пунктиром "--- суффиксные ссылки $\link(v)$.}
\end{center}

\subsubsection{Онлайн-построение и доказательство линейности}

Будем добавлять символы строки $s$ по одному слева направо. Пусть после обработки префикса $s[0 \dots i-1]$ мы знаем вершину $v$, отвечающую наибольшему суффиксу-палиндрому этого префикса (изначально $v = 0$). Добавим символ $c = s[i]$.

Нам нужно найти наибольший суффикс-палиндром новой строки $s[0 \dots i]$. Любой её суффикс-палиндром длины $\ge 2$ (и даже длины $1$, если вспомнить про корень длины $-1$) имеет вид $c + w + c$, где $w$ "--- некоторый суффикс-палиндром префикса $s[0 \dots i-1]$, непосредственно перед которым в строке $s$ стоит буква $c$, то есть выполнено равенство
\[
    s[i - \len(w) - 1] = s[i].
\]
Чтобы длина палиндрома $c + w + c$ была максимальной, нужно взять самый длинный такой суффикс-палиндром $w$ префикса $s[0 \dots i-1]$. Но все суффиксы-палиндромы префикса $s[0 \dots i-1]$ в порядке убывания длины лежат на пути суффиксных ссылок $v$, $\link(v)$, $\link(\link(v))$, $\dots$

Отсюда получаем пошаговый алгоритм добавления символа $s[i]$:
\begin{enumerate}
    \item \textbf{Первый цикл \texttt{while}:} пока $i - \len(v) - 1 < 0$ или $s[i - \len(v) - 1] \ne s[i]$, переходим по суффиксной ссылке $v \la \link(v)$. Заметим, что этот цикл \textbf{всегда завершается}: даже если ни один непустой палиндром и чётный корень $0$ не подошли, мы перейдём в нечётный корень $v = 1$ с длиной $\len(1) = -1$, для которого индекс $i - (-1) - 1$ равен $i$, и условие $s[i] = s[i]$ выполнится автоматически.
    \item Теперь вершина $v$ найдена, и наибольший суффикс-палиндром префикса $s[0 \dots i]$ равен $c + v + c$. Если из $v$ уже есть ребро $\go(v, c)$, то этот палиндром уже встречался раньше: просто делаем $v \la \go(v, c)$ и завершаем шаг.
    \item Если же ребра $\go(v, c)$ ещё нет, создаём новую вершину $u$ с длиной $\len(u) = \len(v) + 2$ и проводим ребро $\go(v, c) = u$. Остаётся найти для новой вершины суффиксную ссылку $\link(u)$:
    \begin{itemize}
        \item Если $v = 1$ (нечётный корень), то новая вершина $u$ отвечает однобуквенному палиндрому длины $1$. Его наибольшим собственным суффиксом-палиндромом является пустая строка, поэтому $\link(u) = 0$.
        \item Если $v \ne 1$, то наибольший собственный суффикс-палиндром строки $u = c + v + c$ непуст и имеет вид $c + v' + c$, где $v'$ "--- самый длинный собственный суффикс-палиндром вершины $v$, перед которым в $s$ стоит буква $c$. Запускаем \textbf{второй цикл \texttt{while}}, стартуя с $v' = \link(v)$, и идём по суффиксным ссылкам $v' \la \link(v')$, пока $s[i - \len(v') - 1] \ne s[i]$. Когда подходящая вершина $v'$ найдена, полагаем $\link(u) = \go(v', c)$ (по лемме~\ref{lem:pal_count} вершина $\go(v', c)$ уже существует в дереве).
    \end{itemize}
    В конце переходим в созданную вершину: $v \la u$.
\end{enumerate}

\begin{lstlisting}[caption={Реализация дерева палиндромов (Eertree) за $O(n)$}]
struct PalindromicTree {
    struct Node {
        vector<int> go;
        int len = 0;
        int link = 0;

        Node(int len = 0, int link = 0) : go(SIGMA, -1), len(len), link(link) {}
    };

    string s;
    vector<Node> t;
    int last; // вершина наибольшего суффикса-палиндрома текущего префикса

    PalindromicTree() {
        t.emplace_back(0, 1);   // вершина 0: чётный корень (len = 0, link = 1)
        t.emplace_back(-1, 1);  // вершина 1: нечётный корень (len = -1, link = 1)
        last = 0;
    }

    int get_link(int v, int i) const {
        while (i - t[v].len - 1 < 0 || s[i - t[v].len - 1] != s[i]) {
            v = t[v].link;
        }
        return v;
    }

    void add_char(char ch) {
        int i = s.size();
        s += ch;
        int c = ch - 'a';

        // 1. Ищем наибольший суффикс-палиндром, расширяемый буквой ch
        int v = get_link(last, i);

        // 2. Если переход уже есть, нового палиндрома не появилось
        if (t[v].go[c] != -1) {
            last = t[v].go[c];
            return;
        }

        // 3. Создаём новую вершину u
        int u = t.size();
        t.emplace_back(t[v].len + 2, 0);
        t[v].go[c] = u;

        if (t[u].len == 1) {
            t[u].link = 0; // ссылка из однобуквенного палиндрома ведёт в чётный корень
        } else {
            int v_prime = get_link(t[v].link, i);
            t[u].link = t[v_prime].go[c];
        }
        last = u;
    }
};
\end{lstlisting}

\begin{theorem}
    Построение дерева палиндромов для строки $s$ длины $n$ работает за время $O(n)$ (с константой $O(|\Sigma|)$ на выделение памяти в узлах, либо за чистое $O(n)$ при хранении переходов списком).
\end{theorem}

\begin{proof}
    По лемме~\ref{lem:pal_count} за всё время работы создаётся не более $n + 2$ вершин. Остаётся оценить суммарное число шагов в двух циклах \texttt{while} внутри вызовов \texttt{get\_link}.

    Проследим за \textbf{позицией левого конца} текущего наибольшего суффикса-палиндрома $v$: после $i$-го шага она равна $L = i - \len(v) + 1$.
    \begin{itemize}
        \item Каждый шаг первого цикла \texttt{while} переводит $v$ в $\link(v)$, при этом длина палиндрома $\len(v)$ строго уменьшается хотя бы на $1$, а значит, левая граница $i - \len(v) + 1$ \textbf{строго возрастает} хотя бы на $1$.
        \item В конце шага мы заменяем $v$ на $\go(v, c)$ с длиной $\len(v) + 2$, а на следующей итерации индекс $i$ увеличивается на $1$. Значит, между итерациями левая граница может уменьшиться не более чем на $1$.
    \end{itemize}
    Так как левая граница всегда лежит в отрезке $[0, n + 1]$ и на каждом из $n$ шагов уменьшается не более чем на $1$, суммарное число её строгих увеличений (то есть шагов первого цикла \texttt{while}) не превосходит $2n$.

    Для второго цикла \texttt{while} рассуждение абсолютно такое же, если следить за левой границей палиндрома $\link(\text{last})$: переход к новой вершине $u$ присваивает $\link(u) = \go(v', c)$, так что длина $\len(\link(\text{last}))$ увеличивается максимум на $2$ относительно $\len(v')$, а каждый шаг второго цикла \texttt{while} строго сдвигает левую границу этого суффикса-палиндрома вправо. Значит, и второй цикл суммарно делает $O(n)$ итераций.
\end{proof}

\subsection{Конечные автоматы и теорема Майхилла--Нероуда}

Перейдём ко второй большой теме "--- \textbf{суффиксному автомату}. Чтобы строго понять, почему эта структура устроена именно так и откуда берётся её минимальность, начнём с классической теории формальных языков.

\begin{definition}
    \textbf{Детерминированным конечным автоматом (ДКА)} называется пятёрка $A = (\Sigma, Q, q_0, T, \delta)$, где:
    \begin{itemize}
        \item $\Sigma$ "--- конечный алфавит;
        \item $Q$ "--- конечное множество состояний автомата;
        \item $q_0 \in Q$ "--- начальное (стартовое) состояние;
        \item $T \subseteq Q$ "--- подмножество терминальных (допускающих) состояний;
        \item $\delta\colon Q \times \Sigma \to Q$ "--- функция переходов (по состоянию и символу однозначно указывающая следующее состояние).
    \end{itemize}
\end{definition}

Функция переходов $\delta$ естественным образом продолжается со символов на произвольные строки $\alpha \in \Sigma^*$ по правилу $\delta(q, \epsilon) = q$ и $\delta(q, \alpha c) = \delta(\delta(q, \alpha), c)$. Говорят, что автомат $A$ \textbf{допускает} строку $\alpha$, если $\delta(q_0, \alpha) \in T$. Множество всех слов, допускаемых автоматом $A$, называется \textbf{языком автомата} и обозначается $L(A) = \{\alpha \in \Sigma^* \mid \delta(q_0, \alpha) \in T\}$.

Пусть нам задан некоторый язык $L \subseteq \Sigma^*$. Как построить ДКА с \emph{минимальным возможным числом состояний}, распознающий язык $L$?

\begin{definition}
    Пусть $L \subseteq \Sigma^*$ и $\alpha \in \Sigma^*$ "--- произвольная строка. \textbf{Правым контекстом} строки $\alpha$ относительно языка $L$ называется множество всех строк $\beta$, при дописывании которых справа к $\alpha$ получается слово из языка $L$:
    \[
        R_L(\alpha) = \{\beta \in \Sigma^* \mid \alpha\beta \in L\}.
    \]
    Две строки $\alpha_1, \alpha_2 \in \Sigma^*$ называются \textbf{контекстно эквивалентными} ($\alpha_1 \sim_L \alpha_2$), если их правые контексты совпадают: $R_L(\alpha_1) = R_L(\alpha_2)$.
\end{definition}

Очевидно, что отношение $\sim_L$ рефлексивно, симметрично и транзитивно, то есть является отношением эквивалентности и разбивает множество всех слов $\Sigma^*$ на непересекающиеся классы эквивалентности $[\alpha]$.

\begin{theorem}[Майхилл--Нероуд]
    Язык $L \subseteq \Sigma^*$ распознаётся некоторым детерминированным конечным автоматом тогда и только тогда, когда отношение эквивалентности $\sim_L$ имеет конечное число классов. При этом минимальное число состояний в ДКА, распознающем язык $L$, в точности равно числу классов эквивалентности отношения $\sim_L$.
\end{theorem}

\begin{proof}[Идея доказательства]
    \textbf{Оценка снизу.} Пусть $A = (\Sigma, Q, q_0, T, \delta)$ "--- произвольный ДКА, распознающий язык $L$. Предположим, что две строки $\alpha_1$ и $\alpha_2$ переводят стартовое состояние $q_0$ в одно и то же состояние $q = \delta(q_0, \alpha_1) = \delta(q_0, \alpha_2)$. Тогда для любого продолжения $\beta \in \Sigma^*$ пути из $q$ по строке $\beta$ совпадают:
    \[
        \delta(q_0, \alpha_1\beta) = \delta(q, \beta) = \delta(q_0, \alpha_2\beta).
    \]
    Значит, $\alpha_1\beta \in L \Lra \alpha_2\beta \in L$, откуда $R_L(\alpha_1) = R_L(\alpha_2)$. Иными словами, строки из разных классов эквивалентности не могут попадать в одно состояние автомата, поэтому $|Q| \ge |\Sigma^* / {\sim_L}|$.

    \textbf{Построение минимального автомата.} Обратно, пусть число классов эквивалентности конечно. Объявим сами классы эквивалентности $[\alpha]$ состояниями автомата! Стартовым состоянием будет класс пустой строки $[\epsilon]$, терминальными "--- те классы $[\alpha]$, для которых $\alpha \in L$ (то есть $\epsilon \in R_L(\alpha)$), а переход по символу $c \in \Sigma$ зададим формулой
    \[
        \delta([\alpha], c) = [\alpha c].
    \]
    Проверим корректность этого определения (независимость от выбора представителя $\alpha$): если $R_L(\alpha_1) = R_L(\alpha_2)$, то для любой строки $\gamma \in \Sigma^*$
    \[
        \gamma \in R_L(\alpha_1 c) \Lra \alpha_1 c \gamma \in L \Lra c\gamma \in R_L(\alpha_1) = R_L(\alpha_2) \Lra \gamma \in R_L(\alpha_2 c),
    \]
    поэтому $R_L(\alpha_1 c) = R_L(\alpha_2 c)$, и переход определён однозначно.
\end{proof}

\subsection{Суффиксный автомат и его структурные свойства}

Применим теперь конструкцию Майхилла--Нероуда к конкретному конечному языку "--- множеству всех суффиксов заданной строки $s$ длины $n$.

\begin{definition}
    \textbf{Суффиксным автоматом} строки $s$ называется минимальный детерминированный конечный автомат, распознающий язык $L = \{s^0, s^1, \dots, s^n\}$ всех суффиксов строки $s$.
\end{definition}

\begin{note}[Об удалении пустого класса]
    Если строка $\alpha$ вообще не является подстрокой строки $s$, то никаким дописыванием символов справа из неё нельзя получить суффикс строки $s$, поэтому $R(\alpha) = \emptyset$. В строгом определении полного ДКА все такие строки образуют один <<мёртвый>> (стоковый) класс эквивалентности $\emptyset$, из которого все переходы ведут в него же. В алгоритмической практике эту стоковую вершину не хранят: вершинами суффиксного автомата считают только непустые классы эквивалентности (состоящие из подстрок строки $s$), а отсутствие перехода из состояния $q$ по букве $c$ просто означает, что переход вёл бы в пустой класс $\emptyset$.
\end{note}

Поймём наглядно, что означает правый контекст $R(\alpha)$ для подстроки $\alpha$ строки $s$. По определению, $\beta \in R(\alpha)$ тогда и только тогда, когда конкатенация $\alpha\beta$ является суффиксом строки $s$. Это значит, что строка $\alpha$ входит в $s$ как подстрока, а $\beta$ "--- это в точности оставшийся <<хвост>> строки $s$ от правого конца этого вхождения $\alpha$ до самого конца строки $s$!

\begin{center}
\begin{tikzpicture}[scale=0.95]
    \draw[thick] (0, 0) rectangle (9, 0.7);
    \draw[thick, fill=NavyBlue!15] (2.5, 0) rectangle (5.0, 0.7) node[pos=0.5, font=\small\bfseries] {$\alpha$};
    \draw[thick, fill=OliveGreen!15] (5.0, 0) rectangle (9, 0.7) node[pos=0.5, font=\small] {$\beta \in R(\alpha)$};
    \node[below, font=\footnotesize] at (0, 0) {$0$};
    \node[below, font=\footnotesize] at (5.0, 0) {$r$};
    \node[below, font=\footnotesize] at (9, 0) {$n$};
    \draw[decorate, decoration={brace, amplitude=6pt}, thick] (2.5, 0.85) -- (9, 0.85) node[midway, above=6pt, font=\small] {суффикс строки $s$};
\end{tikzpicture}
\end{center}

Отсюда видно, что каждое слово $\beta \in R(\alpha)$ однозначно соответствует позиции $r = n - |\beta|$ правого конца очередного вхождения подстроки $\alpha$ в строку $s$. Следовательно, \textbf{две подстроки $\alpha_1$ и $\alpha_2$ лежат в одном состоянии суффиксного автомата ($R(\alpha_1) = R(\alpha_2)$) тогда и только тогда, когда множества правых концов всех их вхождений в строку $s$ полностью совпадают}.

Установим четыре фундаментальные леммы, описывающие устройство классов эквивалентности в суффиксном автомате.

\begin{lemma}[О монотонности правого контекста]\label{lem:sa_mono}
    Если строка $u$ является суффиксом строки $v$, то $R(v) \subseteq R(u)$.
\end{lemma}

\begin{proof}
    Если $u$ "--- суффикс $v$, то $v = w u$ для некоторой строки $w$. Пусть $\beta \in R(v)$. Тогда $v\beta = w u \beta$ является суффиксом строки $s$. Но суффикс суффикса сам является суффиксом, поэтому и $u\beta$ является суффиксом строки $s$, то есть $\beta \in R(u)$.
\end{proof}

\begin{lemma}[О вложенности правых контекстов]\label{lem:sa_nested}
    Пусть $u$ и $v$ "--- две подстроки строки $s$, причём $|u| \le |v|$. Тогда либо их правые контексты не пересекаются ($R(u) \cap R(v) = \emptyset$), либо $u$ является суффиксом $v$ (и тогда по лемме~\ref{lem:sa_mono} выполнено $R(v) \subseteq R(u)$).
\end{lemma}

\begin{proof}
    Предположим, что $R(u) \cap R(v) \ne \emptyset$, и выберем общую строку $\beta \in R(u) \cap R(v)$. Тогда и $u\beta$, и $v\beta$ являются суффиксами одной и той же строки $s$. Это значит, что если отступить от правого конца строки $s$ на $|\beta|$ символов влево, то непосредственно перед хвостом $\beta$ в строке $s$ заканчиваются одновременно и строка $u$, и строка $v$. Так как $|u| \le |v|$, более короткая строка $u$ обязана быть суффиксом более длинной строки $v$!
\end{proof}

\begin{lemma}[Об интервале суффиксов в классе эквивалентности]\label{lem:sa_interval}
    Пусть строка $u$ является суффиксом строки $w$ и $R(u) = R(w)$. Тогда для любого суффикса $v$ строки $w$, промежуточного по длине между $u$ и $w$ ($|u| \le |v| \le |w|$), также выполнено $R(v) = R(u) = R(w)$.
\end{lemma}

\begin{proof}
    Так как $u$ "--- суффикс $v$, а $v$ "--- суффикс $w$, по лемме~\ref{lem:sa_mono} имеем цепочку включений
    \[
        R(w) \subseteq R(v) \subseteq R(u).
    \]
    Но крайние множества в этой цепочке совпадают ($R(w) = R(u)$), значит, совпадает с ними и среднее множество $R(v)$.
\end{proof}

Из лемм~\ref{lem:sa_nested} и~\ref{lem:sa_interval} вытекает полная структурная теорема о классах эквивалентности суффиксного автомата:

\begin{corollary}
    Пусть $C$ "--- произвольное состояние (непустой класс эквивалентности) суффиксного автомата.
    \begin{enumerate}
        \item Любые две строки из класса $C$ имеют общий непустой правый контекст, поэтому по лемме~\ref{lem:sa_nested} более короткая из них всегда является суффиксом более длинной. В частности, в классе $C$ не может быть двух различных строк одинаковой длины!
        \item В классе $C$ есть единственная самая длинная строка $\longest(C)$ и единственная самая короткая строка $\shortest(C)$, а сам класс $C$ состоит \textbf{в точности из всех суффиксов строки $\longest(C)$, имеющих длину от $|\shortest(C)|$ до $|\longest(C)|$ включительно}.
    \end{enumerate}
\end{corollary}

Теперь мы видим, насколько экономно можно описывать состояние $C$: нам не нужно хранить список всех строк, попавших в класс $C$. Достаточно знать его самую длинную строку $\longest(C)$ (её длину мы будем обозначать через $\len(C) = |\longest(C)|$) и то, на какой длине класс $C$ заканчивается! А эту нижнюю границу длины нам подскажет суффиксная ссылка.

\begin{definition}
    Для любого состояния $C$ суффиксного автомата, кроме начального состояния корня $[\epsilon]$, \textbf{суффиксной ссылкой} $\link(C)$ назовём состояние, содержащее самый длинный суффикс строки $\longest(C)$, не лежащий в $C$.
\end{definition}

Какой именно суффикс строки $\longest(C)$ является самым длинным из тех, что уже не попали в $C$? Поскольку в $C$ лежат все суффиксы вплоть до $\shortest(C)$, первый суффикс вне $C$ получается из $\shortest(C)$ отбрасыванием первого символа и имеет длину $|\shortest(C)| - 1$. Более того, в своём новом классе $\link(C)$ этот суффикс обязательно будет \emph{самой длинной} строкой (иначе в $\link(C)$ нашёлся бы ещё более длинный суффикс строки $\longest(C)$, который тогда лежал бы в $C$).

Отсюда получаем замечательную формулу:
\[
    \len(\link(C)) + 1 = |\shortest(C)|,
\]
то есть класс $C$ содержит суффиксы строки $\longest(C)$ всех длин из отрезка $\bigl[\len(\link(C)) + 1,\; \len(C)\bigr]$. Проходя по суффиксным ссылкам $C \to \link(C) \to \link(\link(C)) \to \dots \to [\epsilon]$, мы перебираем все суффиксы строки $\longest(C)$, разбитые на непересекающиеся отрезки длин!

Завершим лекцию ключевым критерием, объясняющим, какие именно подстроки строки $s$ становятся самыми длинными строками ($\longest$) своих классов эквивалентности.

\begin{lemma}[Критерий строки $\longest$]\label{lem:longest_crit}
    Подстрока $\alpha$ строки $s$ является самой длинной строкой в своём классе эквивалентности ($\alpha = \longest([\alpha])$) тогда и только тогда, когда выполняется хотя бы одно из двух условий:
    \begin{enumerate}
        \item $\alpha$ является \textbf{префиксом} всей строки $s$;
        \item существуют две различные буквы $a \ne b$ такие, что обе строки $a\alpha$ и $b\alpha$ являются подстроками строки $s$ (говорят, что строка $\alpha$ \textbf{ветвится влево}).
    \end{enumerate}
\end{lemma}

\begin{proof}
    Условие $\alpha = \longest([\alpha])$ равносильно тому, что ни для какой буквы $c \in \Sigma$ строка $c\alpha$ не лежит в том же классе эквивалентности, что и $\alpha$. Так как $\alpha$ является суффиксом $c\alpha$, по лемме~\ref{lem:sa_mono} всегда $R(c\alpha) \subseteq R(\alpha)$. Значит, $\alpha$ является $\longest$ тогда и только тогда, когда для любой буквы $c \in \Sigma$ включение строгое: $R(c\alpha) \subsetneq R(\alpha)$, то есть у строки $\alpha$ есть хотя бы одно вхождение в $s$, перед которым \emph{не стоит} буква $c$.

    Когда это выполнено для \emph{всех} букв $c \in \Sigma$ одновременно? Возможны ровно два случая:
    \begin{itemize}
        \item Либо у строки $\alpha$ есть вхождение в самом начале строки $s$ (то есть $\alpha$ "--- префикс $s$). Перед этим вхождением вообще нет никакого символа, поэтому при дописывании любой буквы $c$ слева это крайнее левое вхождение теряется, и $R(c\alpha) \subsetneq R(\alpha)$ для всех $c \in \Sigma$.
        \item Либо $\alpha$ не является префиксом $s$, и тогда перед каждым вхождением $\alpha$ в $s$ стоит некоторый символ. Если перед всеми вхождениями $\alpha$ стоит одна и та же буква $a$, то $R(a\alpha) = R(\alpha)$, и $\alpha$ не является $\longest$. Если же перед какими-то двумя вхождениями $\alpha$ стоят разные буквы $a \ne b$ (обе строки $a\alpha$ и $b\alpha$ входят в $s$), то ни одна буква $c \in \Sigma$ не может стоять перед всеми вхождениями $\alpha$ сразу, а значит, $R(c\alpha) \subsetneq R(\alpha)$ для любого $c \in \Sigma$.
    \end{itemize}
\end{proof}
